% Sec. V-A "Reference Map" for main.tex, with the bundle-adjusted and
% TSDF-fused reference (pass 20261015). To use: replace the body of
% \subsection{Reference Map} in main.tex (from "The reference map is built"
% to "...would not be visible to conventional ICP.") with
% % Sec. V-A "Reference Map" for main.tex, with the bundle-adjusted and
% TSDF-fused reference (pass 20261015). To use: replace the body of
% \subsection{Reference Map} in main.tex (from "The reference map is built"
% to "...would not be visible to conventional ICP.") with
% % Sec. V-A "Reference Map" for main.tex, with the bundle-adjusted and
% TSDF-fused reference (pass 20261015). To use: replace the body of
% \subsection{Reference Map} in main.tex (from "The reference map is built"
% to "...would not be visible to conventional ICP.") with
% % Sec. V-A "Reference Map" for main.tex, with the bundle-adjusted and
% TSDF-fused reference (pass 20261015). To use: replace the body of
% \subsection{Reference Map} in main.tex (from "The reference map is built"
% to "...would not be visible to conventional ICP.") with
% \input{sections/refmap_main}. Full method and equations:
% supplementary, supp/refmap_method.tex.
The reference map is built in a dedicated mapping session and held fixed
afterwards; no recorded run contributes to it. A LiDAR-inertial SLAM
solution \cite{glim} gives the initial trajectory. It is refined by plane
bundle adjustment \cite{balm,hba}: the poses of \SI{1}{\second} keyframes
are moved jointly so that every planar patch seen from several keyframes is
as thin as possible, the cost of a patch being the smallest eigenvalue of
the scatter of its points, with the SLAM solution's relative motion and
gravity direction kept as soft priors. The scans are then fused into a
truncated signed distance field \cite{curless1996,vdbfusion} with
\SI{2}{\centi\metre} cells and a truncation of \SI{8}{\centi\metre}, three
standard deviations of the spread with which the scans place a wall; the
map is its zero-level surface, and moving objects are removed by
free-space carving \cite{dufomap}. Fusion matters as much as the poses: the
usual union of points keeps every cell a stray point reached, so on the same
scans and poses its walls are \SI{12.8}{\centi\metre} thick against the
signed distance's \SI{7.8}{\centi\metre}, which averages across a surface
and keeps 98\,\% of the union's surfaces. Each stage is measured by the
thickness of the map's planar surfaces (Table~\ref{tab:refmap_stages}):
from the SLAM solution (time-corrected, below) to the reference, walls go
from \SI{10.7}{\centi\metre} to \SI{7.8}{\centi\metre} and floors from
\SI{2.8}{\centi\metre} to \SI{1.9}{\centi\metre}. Where a wall's remaining
spread comes from is read across the table, by placing the scans that saw
it with each trajectory: the sensor alone, each scan against its own
plane, spreads it over \SI{3.0}{\centi\metre}; a single scan placed in the
map over \SI{4.7}{\centi\metre}; the scans of one second, too
short for drift, over \SI{6.7}{\centi\metre} with either trajectory; the
whole run over \SI{9.0}{\centi\metre} with the reference, against
\SI{12.3}{\centi\metre} with the SLAM solution. Three quarters of the
reference's spread is thus present within one second, and drift adds
\SI{2.3}{\centi\metre} over the \SI{1387}{\second} run, where it added
\SI{5.7}{\centi\metre} before bundle adjustment. Against the reference, the SLAM
trajectory's heading error followed the yaw rate ($r=0.95$), the signature
of a time offset: its pose for time $t$ is the scans' pose at
$t+\SI{50}{\milli\second}$, confirmed independently by the scans agreeing
best with a shift of that size, and not removed by correcting the IMU's own
\SI{11}{\milli\second} lag inside the SLAM. The SLAM labels each pose
with the start of its \SI{100}{\milli\second} sweep, where other LiDAR
odometry uses the middle or the end \cite{molaglim}; ours behaves as the
pose about half a sweep later, so the offset is measured per sensor rather
than assumed, and re-checked on every mapping pass. Corrected on export by
\SI{46}{\milli\second}, which leaves \SI{4}{\milli\second} ($r=0.23$), its
ATE falls from \SI{3.5}{\centi\metre} and $0.72^\circ$ to
\SI{2.7}{\centi\metre} and $0.26^\circ$ (Table~\ref{tab:glim_vs_ref}).
Bundle adjustment started from the SLAM trajectory before and after the
correction reaches the same poses to \SI{0.5}{\centi\metre} and
$0.10^\circ$, and an independent registration of every scan to the map
agrees with it to \SI{1.5}{\centi\metre} and $0.14^\circ$ (ATE): the
result is fixed by the scans, not by its start. The final trajectory closes on the anchor board, seen at
the start and the end of the \SI{1387}{\second} run, within
\SI{14}{\milli\metre}, inside the detection noise.
. Full method and equations:
% supplementary, supp/refmap_method.tex.
The reference map is built in a dedicated mapping session and held fixed
afterwards; no recorded run contributes to it. A LiDAR-inertial SLAM
solution \cite{glim} gives the initial trajectory. It is refined by plane
bundle adjustment \cite{balm,hba}: the poses of \SI{1}{\second} keyframes
are moved jointly so that every planar patch seen from several keyframes is
as thin as possible, the cost of a patch being the smallest eigenvalue of
the scatter of its points, with the SLAM solution's relative motion and
gravity direction kept as soft priors. The scans are then fused into a
truncated signed distance field \cite{curless1996,vdbfusion} with
\SI{2}{\centi\metre} cells and a truncation of \SI{8}{\centi\metre}, three
standard deviations of the spread with which the scans place a wall; the
map is its zero-level surface, and moving objects are removed by
free-space carving \cite{dufomap}. Fusion matters as much as the poses: the
usual union of points keeps every cell a stray point reached, so on the same
scans and poses its walls are \SI{12.8}{\centi\metre} thick against the
signed distance's \SI{7.8}{\centi\metre}, which averages across a surface
and keeps 98\,\% of the union's surfaces. Each stage is measured by the
thickness of the map's planar surfaces (Table~\ref{tab:refmap_stages}):
from the SLAM solution (time-corrected, below) to the reference, walls go
from \SI{10.7}{\centi\metre} to \SI{7.8}{\centi\metre} and floors from
\SI{2.8}{\centi\metre} to \SI{1.9}{\centi\metre}. Where a wall's remaining
spread comes from is read across the table, by placing the scans that saw
it with each trajectory: the sensor alone, each scan against its own
plane, spreads it over \SI{3.0}{\centi\metre}; a single scan placed in the
map over \SI{4.7}{\centi\metre}; the scans of one second, too
short for drift, over \SI{6.7}{\centi\metre} with either trajectory; the
whole run over \SI{9.0}{\centi\metre} with the reference, against
\SI{12.3}{\centi\metre} with the SLAM solution. Three quarters of the
reference's spread is thus present within one second, and drift adds
\SI{2.3}{\centi\metre} over the \SI{1387}{\second} run, where it added
\SI{5.7}{\centi\metre} before bundle adjustment. Against the reference, the SLAM
trajectory's heading error followed the yaw rate ($r=0.95$), the signature
of a time offset: its pose for time $t$ is the scans' pose at
$t+\SI{50}{\milli\second}$, confirmed independently by the scans agreeing
best with a shift of that size, and not removed by correcting the IMU's own
\SI{11}{\milli\second} lag inside the SLAM. The SLAM labels each pose
with the start of its \SI{100}{\milli\second} sweep, where other LiDAR
odometry uses the middle or the end \cite{molaglim}; ours behaves as the
pose about half a sweep later, so the offset is measured per sensor rather
than assumed, and re-checked on every mapping pass. Corrected on export by
\SI{46}{\milli\second}, which leaves \SI{4}{\milli\second} ($r=0.23$), its
ATE falls from \SI{3.5}{\centi\metre} and $0.72^\circ$ to
\SI{2.7}{\centi\metre} and $0.26^\circ$ (Table~\ref{tab:glim_vs_ref}).
Bundle adjustment started from the SLAM trajectory before and after the
correction reaches the same poses to \SI{0.5}{\centi\metre} and
$0.10^\circ$, and an independent registration of every scan to the map
agrees with it to \SI{1.5}{\centi\metre} and $0.14^\circ$ (ATE): the
result is fixed by the scans, not by its start. The final trajectory closes on the anchor board, seen at
the start and the end of the \SI{1387}{\second} run, within
\SI{14}{\milli\metre}, inside the detection noise.
. Full method and equations:
% supplementary, supp/refmap_method.tex.
The reference map is built in a dedicated mapping session and held fixed
afterwards; no recorded run contributes to it. A LiDAR-inertial SLAM
solution \cite{glim} gives the initial trajectory. It is refined by plane
bundle adjustment \cite{balm,hba}: the poses of \SI{1}{\second} keyframes
are moved jointly so that every planar patch seen from several keyframes is
as thin as possible, the cost of a patch being the smallest eigenvalue of
the scatter of its points, with the SLAM solution's relative motion and
gravity direction kept as soft priors. The scans are then fused into a
truncated signed distance field \cite{curless1996,vdbfusion} with
\SI{2}{\centi\metre} cells and a truncation of \SI{8}{\centi\metre}, three
standard deviations of the spread with which the scans place a wall; the
map is its zero-level surface, and moving objects are removed by
free-space carving \cite{dufomap}. Fusion matters as much as the poses: the
usual union of points keeps every cell a stray point reached, so on the same
scans and poses its walls are \SI{12.8}{\centi\metre} thick against the
signed distance's \SI{7.8}{\centi\metre}, which averages across a surface
and keeps 98\,\% of the union's surfaces. Each stage is measured by the
thickness of the map's planar surfaces (Table~\ref{tab:refmap_stages}):
from the SLAM solution (time-corrected, below) to the reference, walls go
from \SI{10.7}{\centi\metre} to \SI{7.8}{\centi\metre} and floors from
\SI{2.8}{\centi\metre} to \SI{1.9}{\centi\metre}. Where a wall's remaining
spread comes from is read across the table, by placing the scans that saw
it with each trajectory: the sensor alone, each scan against its own
plane, spreads it over \SI{3.0}{\centi\metre}; a single scan placed in the
map over \SI{4.7}{\centi\metre}; the scans of one second, too
short for drift, over \SI{6.7}{\centi\metre} with either trajectory; the
whole run over \SI{9.0}{\centi\metre} with the reference, against
\SI{12.3}{\centi\metre} with the SLAM solution. Three quarters of the
reference's spread is thus present within one second, and drift adds
\SI{2.3}{\centi\metre} over the \SI{1387}{\second} run, where it added
\SI{5.7}{\centi\metre} before bundle adjustment. Against the reference, the SLAM
trajectory's heading error followed the yaw rate ($r=0.95$), the signature
of a time offset: its pose for time $t$ is the scans' pose at
$t+\SI{50}{\milli\second}$, confirmed independently by the scans agreeing
best with a shift of that size, and not removed by correcting the IMU's own
\SI{11}{\milli\second} lag inside the SLAM. The SLAM labels each pose
with the start of its \SI{100}{\milli\second} sweep, where other LiDAR
odometry uses the middle or the end \cite{molaglim}; ours behaves as the
pose about half a sweep later, so the offset is measured per sensor rather
than assumed, and re-checked on every mapping pass. Corrected on export by
\SI{46}{\milli\second}, which leaves \SI{4}{\milli\second} ($r=0.23$), its
ATE falls from \SI{3.5}{\centi\metre} and $0.72^\circ$ to
\SI{2.7}{\centi\metre} and $0.26^\circ$ (Table~\ref{tab:glim_vs_ref}).
Bundle adjustment started from the SLAM trajectory before and after the
correction reaches the same poses to \SI{0.5}{\centi\metre} and
$0.10^\circ$, and an independent registration of every scan to the map
agrees with it to \SI{1.5}{\centi\metre} and $0.14^\circ$ (ATE): the
result is fixed by the scans, not by its start. The final trajectory closes on the anchor board, seen at
the start and the end of the \SI{1387}{\second} run, within
\SI{14}{\milli\metre}, inside the detection noise.
. Full method and equations:
% supplementary, supp/refmap_method.tex.
The reference map is built in a dedicated mapping session and held fixed
afterwards; no recorded run contributes to it. A LiDAR-inertial SLAM
solution \cite{glim} gives the initial trajectory. It is refined by plane
bundle adjustment \cite{balm,hba}: the poses of \SI{1}{\second} keyframes
are moved jointly so that every planar patch seen from several keyframes is
as thin as possible, the cost of a patch being the smallest eigenvalue of
the scatter of its points, with the SLAM solution's relative motion and
gravity direction kept as soft priors. The scans are then fused into a
truncated signed distance field \cite{curless1996,vdbfusion} with
\SI{2}{\centi\metre} cells and a truncation of \SI{8}{\centi\metre}, three
standard deviations of the spread with which the scans place a wall; the
map is its zero-level surface, and moving objects are removed by
free-space carving \cite{dufomap}. Fusion matters as much as the poses: the
usual union of points keeps every cell a stray point reached, so on the same
scans and poses its walls are \SI{12.8}{\centi\metre} thick against the
signed distance's \SI{7.8}{\centi\metre}, which averages across a surface
and keeps 98\,\% of the union's surfaces. Each stage is measured by the
thickness of the map's planar surfaces (Table~\ref{tab:refmap_stages}):
from the SLAM solution (time-corrected, below) to the reference, walls go
from \SI{10.7}{\centi\metre} to \SI{7.8}{\centi\metre} and floors from
\SI{2.8}{\centi\metre} to \SI{1.9}{\centi\metre}. Where a wall's remaining
spread comes from is read across the table, by placing the scans that saw
it with each trajectory: the sensor alone, each scan against its own
plane, spreads it over \SI{3.0}{\centi\metre}; a single scan placed in the
map over \SI{4.7}{\centi\metre}; the scans of one second, too
short for drift, over \SI{6.7}{\centi\metre} with either trajectory; the
whole run over \SI{9.0}{\centi\metre} with the reference, against
\SI{12.3}{\centi\metre} with the SLAM solution. Three quarters of the
reference's spread is thus present within one second, and drift adds
\SI{2.3}{\centi\metre} over the \SI{1387}{\second} run, where it added
\SI{5.7}{\centi\metre} before bundle adjustment. Against the reference, the SLAM
trajectory's heading error followed the yaw rate ($r=0.95$), the signature
of a time offset: its pose for time $t$ is the scans' pose at
$t+\SI{50}{\milli\second}$, confirmed independently by the scans agreeing
best with a shift of that size, and not removed by correcting the IMU's own
\SI{11}{\milli\second} lag inside the SLAM. The SLAM labels each pose
with the start of its \SI{100}{\milli\second} sweep, where other LiDAR
odometry uses the middle or the end \cite{molaglim}; ours behaves as the
pose about half a sweep later, so the offset is measured per sensor rather
than assumed, and re-checked on every mapping pass. Corrected on export by
\SI{46}{\milli\second}, which leaves \SI{4}{\milli\second} ($r=0.23$), its
ATE falls from \SI{3.5}{\centi\metre} and $0.72^\circ$ to
\SI{2.7}{\centi\metre} and $0.26^\circ$ (Table~\ref{tab:glim_vs_ref}).
Bundle adjustment started from the SLAM trajectory before and after the
correction reaches the same poses to \SI{0.5}{\centi\metre} and
$0.10^\circ$, and an independent registration of every scan to the map
agrees with it to \SI{1.5}{\centi\metre} and $0.14^\circ$ (ATE): the
result is fixed by the scans, not by its start. The final trajectory closes on the anchor board, seen at
the start and the end of the \SI{1387}{\second} run, within
\SI{14}{\milli\metre}, inside the detection noise.
